\subsection{维数的半连续性}

引理 5.--- 设 A 是一个环，r 是其根基，\(\mathbf{R} = \bigoplus_{i \in \mathbf{Z}} \mathbf{R}_i\) 是一个分次 A-代数，\(\mathbf{M} = \bigoplus_{i \in \mathbf{Z}} \mathbf{M}_i\) 是一个分次 R-模。假设每个 \(\mathbf{M}_i\) 都是有限型的 A-模

并且 M/rM 是有限生成的 R/rR-模。则 M 是有限生成的 R-模。设 \(m_{1}, \ldots, m_{n}\) 是 M 中的齐次元素，它们在 M/rM 中的像生成 R/rR-模 M/rM。令 N 为由 \(\{m_{1}, \ldots, m_{n}\}\) 生成的 M 的（分次）子 R-模。对于每个 \(i \in \mathbf{Z}\)，有 \(M_{i} = N_{i} + rM_{i}\)，因而 \(M_{i} = N_{i}\)（II, §3, 第 2 号，命题 4）；故 M = N。

引理 6. --- 设 \(\rho: \mathbf{B} \to \mathbf{C}\) 是环同态，S 是 B 的一个乘法子集。假设 C 是有限生成的 B-代数，且 \(\mathbf{S}^{-1}\mathbf{C}\) 是有限的 \(\mathbf{S}^{-1}\mathbf{B}\)-代数。则存在 \(f \in \mathbf{S}\)，使得 \(\mathbf{C}_f\) 是有限的 \(\mathbf{B}_f\)-代数。

设 X 是 B-代数 C 的一个有限生成集。对于每个 \(x \in X\)，x 在 \(S^{-1}C\) 中的像在 \(S^{-1}B\) 上整，因此存在整数 \(n(x) \geqslant 0\)、元素 \(b_{1}(x), \ldots, b_{n(x)}(x) \in B\) 以及元素 \(f(x) \in S\)，使得

\[
f (x) x ^ {n (x)} + b _ {1} (x) x ^ {n (x) - 1} + \dots + b _ {n (x)} = 0.
\]

置 \(f = \prod_{x \in X} f(x)\)；对于每个 \(x\)，\(X\) 中的该元素在 \(C_f\) 中的像都在 \(B_f\) 上整，故 \(C_f\) 是有限的 \(B_f\)-代数（V, § 1, 第 1 号，命题 4）。

命题 7. — 假设 H 是有限型的 \(H_{0}\)-代数。则函数 \(\mathfrak{p} \mapsto \dim(\mathrm{H} \otimes_{\mathrm{H}_{0}} \kappa(\mathfrak{p}))\) 在 \(\operatorname{Spec}(\mathrm{H}_{0})\) 上是上半连续的。

由于 H 作为 \(H_{0}\)-代数是有限生成的，每个 \(H_{i}\) 都是有限生成的 \(H_{0}\)-模（III, § 1, 第 2 号，命题 1 的推论），并且 H 作为 \(H_{0}\)-代数由 \(H_{0} \oplus H_{1} \oplus \cdots \oplus H_{r}\) 生成，其中 \(r \geqslant 0\) 是适当的整数。取 \(\mathfrak{p} \in \operatorname{Spec}(H_{0})\)，置 \(\dim(H \otimes_{H_{0}} \kappa(\mathfrak{p})) = n \geqslant 0\)。根据命题 6 的推论 1，存在 H 中的元素 \(a_{1}, ..., a_{n}\)，它们都是同一次数 \(d > 0\) 的齐次元素，使得 \(\kappa(\mathfrak{p})\)-同态 \(\overline{\varphi}: \kappa(\mathfrak{p})[X_{1}, ..., X_{n}] \to H \otimes_{H_{0}} \kappa(\mathfrak{p})\) 将 \(X_{i}\) 送到 \(a_{i} \otimes 1\)（\(1 \leqslant i \leqslant n\)），并使 \(H \otimes_{H_{0}} \kappa(\mathfrak{p})\) 成为有限的 \(\kappa(\mathfrak{p})[X_{1}, ..., X_{n}]\)-代数。令 \(\varphi\) 为一个 \(H_{0}\)-同态，它从 \(H_{0}[X_{1}, ..., X_{n}] = R\) 到 H，并将 \(X_{i}\) 送到 \(a_{i}\)，其中 \(1 \leqslant i \leqslant n\)。若对每个 \(m \in \mathbf{Z}\) 置 \(H_{m}' = \sum_{(m-1)d < i \leqslant md} H_{i}\)，则在 H 上得到一个 \(\mathbf{Z}\)-型分次，且与由 \(\varphi\) 给出的 R-模结构相容。每个 \(H_{m}'\) 都是有限型的 \(H_{0}\)-模。根据引理 5，\(H_{\mathfrak{p}}\) 是有限生成的 \(R_{\mathfrak{p}}\)-模。根据引理 6，存在 \(f \in H_{0} - \mathfrak{p}\)，使得 \(H_{f}\) 是有限型的 \(R_{f}\)-模。对于每个 \(\mathfrak{q} \in \operatorname{Spec}(H_{0})_{f}\)，\(H \otimes_{H_{0}} \kappa(\mathfrak{q})\) 是有限的 \(\kappa(\mathfrak{q})[X_{1}, ..., X_{n}]\)-代数，故 \(\dim(H \otimes_{H_{0}} \kappa(\mathfrak{q})) \leqslant n\)（\S 2，第 3 号，定理 1），这就完成了证明。

注记。--- 1) * 在代数几何中，命题 7 蕴含：代数簇之间的射影态射的纤维维数是上半连续的。 *

\begin{enumerate}
\def\labelenumi{\arabic{enumi})}
\setcounter{enumi}{1}
\tightlist
\item
  我们稍后将看到，若 \(\rho:A\to B\) 是环同态，使 \(B\) 成为有限型的 \(A\)-代数，则函数 \(\mathfrak{q}\mapsto\dim_{\mathfrak{q}}(B\otimes_{A}\kappa(\rho^{-1}(q)))\) 在 \(\operatorname{Spec}(B)\) 上是上半连续的（参见 p. 101，练习 10）。
\end{enumerate}

